\section{Preprocesado}
\subsection{Descarga de datos y modelos}

Para garantizar la reproducibilidad y facilitar el despliegue del proyecto, se ha implementado un sistema automatizado de gestión de recursos siguiendo principios de diseño \texttt{SOLID} y \texttt{Clean Code}. Este sistema se encarga de la descarga y organización tanto del dataset como de los modelos de lenguaje.

\subsubsection{Arquitectura del Sistema de Ingestión}

La arquitectura se ha desacoplado para que cada dominio (datasets y modelos) gestione sus propios recursos:

\begin{itemize}
    \item \textbf{Gestión de Modelos (\texttt{.models/}):} Se ha centralizado la descarga de modelos en un directorio oculto en la raíz del proyecto, manteniendo la lógica desacoplada en \texttt{app/models/}.
    
\lstinputlisting[language=Python, firstline=3, lastline=19, caption=Lógica de descarga de modelos en app/models/downloader.py]{app/models/downloader.py}

    \item \textbf{Gestión de Datasets (\texttt{app/dataset/}):} Se encarga de la clonación del repositorio de datos mediante Git, manteniendo una interfaz abstracta para futuros cambios.

    \item \textbf{Configuración Dinámica:} Se utilizan clases especializadas para calcular las rutas de forma robusta, asegurando que el proyecto funcione independientemente del directorio de ejecución.
\end{itemize}

\subsubsection{Flujo de Ejecución}

El punto de entrada principal coordina ambos servicios, inyectando los descargadores específicos para cada tipo de recurso.

\lstinputlisting[language=Python, firstline=7, lastline=35, caption=Orquestación del setup en app/main.py]{app/main.py}

Este enfoque modular permite que, por ejemplo, el sistema de descarga de modelos pueda ser modificado o testeado sin afectar a la lógica de gestión de los datasets.

\subsection{Tokenización}
\subsection{Truncado}
\subsection{Padding}
\subsection{Manejo Secuencias Largas}
